\documentclass[12pt]{article}
\usepackage[utf8]{inputenc}
\usepackage[T1]{fontenc}
\usepackage[margin=1in]{geometry}
\usepackage{setspace}
\usepackage{xcolor}
\usepackage[colorlinks=true, urlcolor=blue!60!black, linkcolor=blue!60!black]{hyperref}
\setlength{\parindent}{2em}
\setlength{\parskip}{1.25em}
\renewcommand{\baselinestretch}{1.0}
\begin{document}
\begin{center}
{\LARGE \textit{Baker v.\ Carr}, 369 U.S.\ 186 (1962)\par}
\vspace{0.8em}
{\large Case Brief\par}
\end{center}
\vspace{1em}

\section{Memory Jogger}

In \href{https://supreme.justia.com/cases/federal/us/369/186/}{\textit{Baker v. Carr}}, the Supreme Court held that Fourteenth Amendment Equal Protection challenges to malapportioned state legislative districts are justiciable in federal court and articulated the modern six-factor test governing the political question doctrine.

\section{Detailed Case Facts}

\textbf{Factual Background.} Tennessee's General Assembly consisted of a 33-member Senate and a 99-member House of Representatives apportioned among the State's 95 counties. Under Article II, \S\S~4--6 of the Tennessee Constitution, the legislature was required to enumerate qualified voters every ten years starting in 1871 and reapportion legislative seats among counties or districts according to the number of qualified voters in each. Although the General Assembly complied from 1871 through 1901, it enacted no general reapportionment after passing the Apportionment Act of 1901 (Tenn. Pub. Acts 1901, ch. 122, codified at Tenn. Code Ann. \S\S~3-101 to 3-107). Over the next six decades, every reapportionment bill introduced in either chamber failed.

Between 1901 and 1961, Tennessee's population grew from 2,020,616 (with 487,380 eligible voters) to 3,567,089 (with 2,092,891 eligible voters) according to the 1960 Federal Census, accompanied by dramatic urbanization and migration from rural counties into metropolitan centers such as Memphis (Shelby County), Nashville (Davidson County), Chattanooga (Hamilton County), and Knoxville (Knox County). Because the 1901 map remained frozen, severe voting disparities emerged: approximately 37\% of Tennessee's voters elected 20 of the 33 Senators, and 40\% elected 63 of the 99 House members. Justice Douglas noted allegations that a single vote in rural Moore County carried the weight of 19 votes in urban Hamilton County, while Justice Clark demonstrated from 1950 voting-population data that Moore County (2,340 voters) enjoyed the same total representation as Rutherford County (25,316 voters), and that arbitrary disparities persisted even among counties of comparable population.

\textbf{Parties and Statutory/Constitutional Provisions.} Charles W. Baker (a qualified voter in Shelby County) and other voters from Davidson, Hamilton, Knox, Montgomery, and Shelby Counties---joined as intervenors by the Mayor of Nashville and the Cities of Chattanooga and Knoxville---sued Tennessee Secretary of State Joe C. Carr, Attorney General George F. McCanless, Coordinator of Elections Jerry McDonald, and members of the State Board of Elections under 42 U.S.C. \S\S~1983 and 1988, invoking federal subject-matter jurisdiction under 28 U.S.C. \S~1343(3). Plaintiffs alleged that the 1901 Act was arbitrary even when enacted and had become ``unconstitutional and obsolete,'' inflicting a ``debasement of their votes'' in violation of the Equal Protection Clause of the Fourteenth Amendment. They sought a declaratory judgment invalidating the 1901 Act, an injunction barring further elections under it, and either a judicially decreed reapportionment or at-large elections pending valid legislative action. Defendants invoked the political question doctrine and argued that any challenge to state legislative structure sounded exclusively under the nonjusticiable Guaranty Clause of Article IV, \S~4 (``The United States shall guarantee to every State in this Union a Republican Form of Government'').

\textbf{Historical Context.} By 1962, rapid mid-twentieth-century suburbanization and urbanization had left a majority of state legislatures dominated by rural minorities entrenched under decades-old districting maps. Tennessee offered no popular initiative or referendum, and constitutional conventions could be called only by the malapportioned General Assembly itself. Moreover, the Tennessee Supreme Court had declined to grant relief in \textit{Kidd v. McCanless}, 200 Tenn. 273, 292 S.W.2d 40 (1956), leaving urban voters with no effective state political or judicial remedy.

\section{Procedural History}

\textbf{Three-Judge District Court.} A single federal district judge determined that plaintiffs' claims were not obviously meritless and convened a statutory three-judge court under 28 U.S.C. \S~2281 in the United States District Court for the Middle District of Tennessee (\textit{Baker v. Carr}, 175 F. Supp. 649 (M.D. Tenn. 1959)). In \href{https://law.justia.com/cases/federal/district-courts/FSupp/179/824/1522687/}{\textit{Baker v. Carr}, 179 F. Supp. 824 (M.D. Tenn. 1959)}, the three-judge panel issued a per curiam opinion dismissing the complaint on two grounds: (1) lack of subject-matter jurisdiction and (2) failure to state a claim upon which relief could be granted. Remarkably, the panel agreed that Tennessee's legislature was ``guilty of a clear violation of the state constitution and of the rights of the plaintiffs,'' yet concluded that \href{https://supreme.justia.com/cases/federal/us/328/549/}{\textit{Colegrove v. Green}, 328 U.S. 549 (1946)}, and subsequent per curiam dismissals foreclosed federal judicial intervention in legislative apportionment.

\textbf{Supreme Court Review and Disposition.} Plaintiffs took a direct appeal to the U.S. Supreme Court pursuant to 28 U.S.C. \S~1253. The Court noted probable jurisdiction (364 U.S. 898 (1960)), heard oral argument on April 19--20, 1961, ordered reargument on May 1, 1961 (366 U.S. 907), and heard reargument on October 9, 1961, with Solicitor General Archibald Cox arguing as amicus curiae in support of reversal (\href{https://www.oyez.org/cases/1960/6}{Oyez summary}). On March 26, 1962, the Court \textbf{reversed and remanded} the case to the District Court for trial and further proceedings. (On remand, the District Court evaluated subsequent legislative enactments in \href{https://law.justia.com/cases/federal/district-courts/FSupp/222/684/2009237/}{\textit{Baker v. Carr}, 222 F. Supp. 684 (M.D. Tenn. 1963)}.)

\section{Judicial Votes}

The Court decided the case by a \textbf{6--2} vote (\href{https://supreme.justia.com/cases/federal/us/369/186/}{369 U.S. at 187, 237}):

\begin{itemize}
  \item \textbf{Majority (6):} Justice William J. Brennan, Jr. delivered the opinion of the Court, joined by Chief Justice Earl Warren and Justices Hugo L. Black, William O. Douglas, Tom C. Clark, and Potter Stewart.
  \item \textbf{Concurrences (3):} Justice William O. Douglas filed a concurring opinion (joining the opinion of the Court); Justice Tom C. Clark filed a concurring opinion (joining the opinion of the Court); Justice Potter Stewart filed a concurring opinion (joining the opinion of the Court).
  \item \textbf{Dissents (2):} Justice Felix Frankfurter filed a dissenting opinion, joined by Justice John Marshall Harlan II. Justice John Marshall Harlan II also filed a separate dissenting opinion.
  \item \textbf{Not Participating (1):} Justice Charles E. Whittaker did not participate in the decision of the case (and retired shortly thereafter).
\end{itemize}

\section{Holding}

\textbf{Judgment.} The judgment of the three-judge District Court dismissing the complaint is reversed, and the cause is remanded for trial and further proceedings.

\textbf{Rule Announced.} The Court issued three threshold holdings (\href{https://supreme.justia.com/cases/federal/us/369/186/}{369 U.S. at 197--98}): (1) a federal district court possesses \textbf{subject-matter jurisdiction} under Article III, \S~2 and 28 U.S.C. \S~1343(3) over a non-frivolous claim that a state legislative apportionment statute violates the Fourteenth Amendment; (2) qualified voters alleging arbitrary vote dilution under such a statute have \textbf{standing} because they assert a concrete personal stake in maintaining the effectiveness of their votes; and (3) a Fourteenth Amendment Equal Protection challenge to state legislative malapportionment presents a \textbf{justiciable cause of action} rather than a nonjusticiable political question. In articulating the political question doctrine, Justice Brennan synthesized six independent formulations, declaring that a case should not be dismissed as a political question unless one of those six attributes is ``inextricable from the case at bar'' (\href{https://supreme.justia.com/cases/federal/us/369/186/}{369 U.S. at 217}).

\textbf{Holding vs. Dicta.} \textit{Baker}'s strict holding is purely jurisdictional and procedural: federal courts may hear Equal Protection challenges to state legislative malapportionment. The Court expressly declined to decide the merits of Tennessee's 1901 Act, the substantive constitutional standard governing apportionment (later resolved as ``one person, one vote'' in \href{https://supreme.justia.com/cases/federal/us/376/1/}{\textit{Wesberry v. Sanders}, 376 U.S. 1 (1964)}, and \href{https://supreme.justia.com/cases/federal/us/377/533/}{\textit{Reynolds v. Sims}, 377 U.S. 533 (1964)}), or the appropriate judicial remedy. Strictly speaking, Justice Brennan's sweeping taxonomy of political question precedents across foreign relations, war duration, constitutional amendments, Indian tribes, and the Guaranty Clause was analytical synthesis supporting the distinction from \href{https://supreme.justia.com/cases/federal/us/48/1/}{\textit{Luther v. Borden}, 48 U.S. 1 (1849)}, though the six-factor test has since been treated by the Court as the canonical doctrinal framework (see, e.g., \href{https://supreme.justia.com/cases/federal/us/506/224/}{\textit{Nixon v. United States}, 506 U.S. 224 (1993)}; \href{https://supreme.justia.com/cases/federal/us/566/189/}{\textit{Zivotofsky v. Clinton}, 566 U.S. 189 (2012)}).

\section{Analysis of Opinions}

\subsection{Majority Opinion (Justice Brennan)}

Justice Brennan carefully separated three concepts that the District Court had conflated: subject-matter jurisdiction, standing, and justiciability. First, under \textit{Bell v. Hood}, 327 U.S. 678 (1946), a federal claim confers Article III and statutory jurisdiction under 28 U.S.C. \S~1343(3) unless it is ``so attenuated and unsubstantial as to be absolutely devoid of merit'' or ``frivolous.'' Re-reading \href{https://supreme.justia.com/cases/federal/us/328/549/}{\textit{Colegrove v. Green}}, Brennan noted that four of the seven participating Justices (Black, Douglas, Murphy, and Rutledge) had agreed that subject-matter jurisdiction existed, and that Justice Rutledge's decisive fourth vote to affirm rested solely on equitable discretion given the proximity of the impending election.

Second, on standing, Brennan framed the classic test: whether appellants alleged ``such a personal stake in the outcome of the controversy as to assure that concrete adverseness which sharpens the presentation of issues.'' Because plaintiffs asserted an individual injury---arbitrary impairment of the effectiveness of their votes---they did not merely raise a generalized citizen grievance.

Third, on justiciability, Brennan reoriented the political question doctrine around \textbf{horizontal separation of powers}: ``it is the relationship between the judiciary and the coordinate branches of the Federal Government, and not the federal judiciary's relationship to the States, which gives rise to the `political question.''' Surveying prior cases in foreign affairs, duration of hostilities, enactment validity, Indian tribal status, and the Guaranty Clause, Brennan distilled six factors: (1) a textually demonstrable constitutional commitment to a coordinate political department; (2) a lack of judicially discoverable and manageable standards; (3) the impossibility of deciding without an initial nonjudicial policy determination; (4) the impossibility of independent resolution without expressing lack of respect due coordinate branches; (5) an unusual need for unquestioning adherence to a political decision already made; or (6) the potentiality of embarrassment from multifarious pronouncements by various departments.

Applying this framework, Brennan concluded that Guaranty Clause claims under \href{https://supreme.justia.com/cases/federal/us/48/1/}{\textit{Luther v. Borden}} are nonjusticiable only because they implicate those political-question elements---not because they touch state governmental structure. By contrast, plaintiffs' Equal Protection claim involved no coordinate federal branch, risked no foreign embarrassment, and invoked ``well developed and familiar'' judicial standards capable of identifying arbitrary and capricious state discrimination, just as the Court had enforced the Fifteenth Amendment against municipal gerrymandering in \href{https://supreme.justia.com/cases/federal/us/364/339/}{\textit{Gomillion v. Lightfoot}, 364 U.S. 339 (1960)}.

\subsection{Concurring Opinion (Justice Douglas)}

Joining the Court's opinion, Justice Douglas emphasized that voting rights are foundational to representative government and have long been protected by federal courts against race discrimination, ballot-box stuffing, and false tallies. While acknowledging that the Equal Protection Clause permits some weighting and bars only ``invidious discrimination,'' Douglas argued that if plaintiffs' allegations (such as a 19-to-1 voting disparity between Moore and Hamilton Counties) were proved, a clear constitutional violation was established. In footnotes, Douglas went further than Brennan, questioning whether \textit{Luther v. Borden}'s broad Guaranty Clause abstention remained defensible and noting that entrenched legislatures make nonjudicial relief illusory.

\subsection{Concurring Opinion (Justice Clark)}

Justice Clark joined the Court's opinion and acquiesced in the remand, but criticized the majority for refusing to decide the merits and give the District Court guidance when the record was undisputed. Applying rational-basis Equal Protection scrutiny under \textit{MacDougall v. Green}, 335 U.S. 281 (1948), Clark examined county-by-county data in an appended Table I and concluded that Tennessee's apportionment was ``a crazy quilt without rational basis,'' discriminating horizontally among rural counties and among urban counties alike. Crucially, Clark stressed that he would not have intervened had Tennessee voters possessed any practical political remedy; because Tennessee lacked initiative and referendum and the legislature had locked itself into power, the electorate was ``caught up in a legislative strait jacket,'' leaving federal courts as the sole avenue of relief.

\subsection{Concurring Opinion (Justice Stewart)}

Justice Stewart wrote a brief, restraining concurrence to clarify what the Court did \textit{not} decide. Stressing that the Court decided only jurisdiction, standing, and justiciability, Stewart insisted that the opinion neither required ``approximate equality'' of every vote nor barred a State, ``acting not irrationally,'' from choosing a legislative structure suited to its people's interests and customs. Because Tennessee had not yet even filed an answer, the merits belonged to the District Court in the first instance.

\subsection{Dissenting Opinion (Justice Frankfurter, joined by Justice Harlan)}

In his final major opinion before suffering a stroke, Justice Frankfurter delivered a passionate defense of judicial restraint. Calling the decision ``a massive repudiation of the experience of our whole past,'' Frankfurter argued that \textit{Baker} overruled a uniform line of cases from \textit{Colegrove v. Green} through \textit{Kidd v. McCanless}. He warned that the Court's authority---possessing ``neither the purse nor the sword''---depends on sustained public confidence nourished by detachment from ``political entanglements.''

Substantively, Frankfurter contended that plaintiffs' Equal Protection claim was merely a ``Guaranty Clause claim masquerading under a different label.'' Deciding how to allocate legislative representation required choosing among competing theories of political philosophy (numerical equality, geographic balance, urban-rural representation, historic subdivisions), for which no judicial standards existed. Catapulting lower courts into this ``mathematical quagmire'' without remedies or standards would invite legislatures to ``play ducks and drakes with the judiciary.'' Relief, he insisted, belonged not to courts but to ``an informed, civically militant electorate.''

\subsection{Dissenting Opinion (Justice Harlan)}

Justice Harlan wrote separately to argue that the complaint should be dismissed not merely for nonjusticiability, but for failure to state a claim on the merits under the Equal Protection Clause. Surveying American and British political history, Harlan maintained that nothing in the Fourteenth Amendment requires state legislatures ``to reflect with approximate equality the voice of every voter.'' A State's choice to weight geography or balance rural and urban interests was a classic legislative judgment entitled to deference, and Tennessee's 1901 Act was not demonstrably irrational once one accounted for floterial districts and historical county representation. Harlan lamented that the majority treated the Court as ``the last refuge for the correction of all inequality or injustice.''

\section{Examples: Future Applications}

\begin{enumerate}
  \item \textbf{Same-Side (Justiciable --- Passport Designation Statute):} Congress enacts a statute directing the Secretary of State, upon request, to record ``Israel'' as the place of birth on passports of U.S. citizens born in Jerusalem. The Executive refuses, claiming exclusive recognition power. \textit{Outcome:} Justiciable under \textit{Baker}. As the Court held in \href{https://supreme.justia.com/cases/federal/us/566/189/}{\textit{Zivotofsky v. Clinton}, 566 U.S. 189 (2012)}, the plaintiff does not ask the court to set foreign policy regarding Jerusalem's sovereignty, but merely to perform the familiar judicial task of determining whether a federal statute is constitutional or impermissibly infringes executive power.

  \item \textbf{Same-Side (Justiciable --- State Legislative Malapportionment on the Merits):} Voters in Alabama challenge a state constitutional scheme guaranteeing at least one representative per county, producing a 41-to-1 population disparity between Jefferson County and a rural county. \textit{Outcome:} Justiciable (and unconstitutional) under \textit{Baker}. Following \textit{Baker}'s holding that Equal Protection vote-dilution claims against state legislatures present no separation-of-powers bar and are amenable to judicial standards, the Court struck down the scheme in \href{https://supreme.justia.com/cases/federal/us/377/533/}{\textit{Reynolds v. Sims}, 377 U.S. 533 (1964)}.

  \item \textbf{Opposite-Side (Nonjusticiable --- Senate Impeachment Trial Procedures):} A federal judge impeached by the House challenges his Senate conviction, arguing that the Senate's use of a twelve-member committee to receive evidence rather than a full-floor hearing violates Article I, \S~3, cl. 6's requirement that the Senate ``try'' all impeachments. \textit{Outcome:} Nonjusticiable under \textit{Baker} Factors 1 and 2. As held in \href{https://supreme.justia.com/cases/federal/us/506/224/}{\textit{Nixon v. United States}, 506 U.S. 224 (1993)}, the word ``sole'' in ``sole Power to try all Impeachments'' is a textually demonstrable commitment to the Senate, and ``try'' lacks judicially manageable standards; moreover, judicial review would risk constitutional chaos and lack of finality.

  \item \textbf{Opposite-Side (Nonjusticiable --- Partisan Gerrymandering):} Voters challenge a state's congressional map drawn with equal-population districts, alleging that extreme partisan ``cracking and packing'' dilutes their political influence in violation of the Equal Protection Clause and First Amendment. \textit{Outcome:} Nonjusticiable under \textit{Baker} Factor 2. In \href{https://supreme.justia.com/cases/federal/us/588/18-422/}{\textit{Rucho v. Common Cause}, 588 U.S. 684 (2019)}, the Court distinguished \textit{Baker} and \textit{Reynolds} because ``one person, one vote'' offers a clear numerical rule, whereas determining how much partisan consideration is ``too much'' lacks a limited, precise, judicially discernible and manageable standard.

  \item \textbf{Fence-Sitter (Borderline --- War Powers Resolution Challenge):} Members of Congress sue the President to enjoin ongoing overseas military airstrikes lasting beyond 60 days without congressional authorization under the War Powers Resolution. \textit{Why It Is Close:} Under \textit{Baker}, courts must separate statutory interpretation from military discretion. Challengers argue---like in \textit{Zivotofsky}---that construing a federal statute and Article I's Declare War Clause is a core judicial duty with manageable legal standards. The Executive counters that deciding whether ``hostilities'' exist and whether to withdraw troops implicates foreign-relations commitment (Factor 1), lack of manageable standards (Factor 2), and the risk of multifarious pronouncements abroad (Factor 6), leading lower federal courts frequently to dismiss such suits on standing, ripeness, or political-question grounds (see \href{https://www.everycrsreport.com/reports/R43834.html}{CRS Report R43834}).
\end{enumerate}

\section{Critique}

\textbf{Logical Weaknesses and Doctrinal Tensions.} Scholars have long noted that \textit{Baker}'s six-factor test is both over-inclusive and under-theorized. Justice Brennan purported merely to summarize past decisions, yet he quietly transformed the political question doctrine from a pragmatic mix of federalism, equity, and institutional deference into a rigid, federal-separation-of-powers-only inquiry. As Professor Tara Leigh Grove has documented in \href{https://scholarship.law.wm.edu/cgi/viewcontent.cgi?article=2834&context=facpubs}{\textit{The Lost History of the Political Question Doctrine}}, \textit{Baker} simultaneously asserted judicial supremacy---claiming that determining whether another branch has constitutional authority ``is a responsibility of this Court as ultimate interpreter of the Constitution''---while narrowing deference to coordinate branches. Moreover, Brennan's assertion that ``well developed and familiar'' Equal Protection standards were readily available to judge legislative apportionment was largely ipse dixit: \textit{Baker} offered lower courts no substantive standard at all, forcing the Court two years later in \textit{Reynolds v. Sims} to adopt a bright-line ``one person, one vote'' rule that Justice Stewart's concurrence had explicitly disclaimed.

\textbf{Originalist and Textualist Perspectives.} From an originalist and textualist standpoint, \textit{Baker} (and its progeny) sits uneasily with the text and history of the Fourteenth Amendment. As Justice Harlan stressed in dissent---and as originalist critics echo---Section 2 of the Fourteenth Amendment explicitly contemplated that States might restrict or structure the franchise (imposing reduced congressional representation as the penalty for denying the vote to adult males), strongly suggesting that Section 1's Equal Protection Clause was not originally understood to regulate state legislative districting. Originalists also question the prudential prongs of \textit{Baker}'s test (Factors 4, 5, and 6), arguing that courts have a duty to decide cases within their Article III jurisdiction unless the constitutional text itself commits the question to another branch (Factor 1) or supplies no legal rule (Factor 2)---a narrowing trend visible in \href{https://supreme.justia.com/cases/federal/us/566/189/}{\textit{Zivotofsky v. Clinton}}.

\textbf{Progressive and Representation-Reinforcing Perspectives.} Progressive and process-based theorists, most famously John Hart Ely, celebrate \textit{Baker} as a triumph of ``representation-reinforcing'' judicial review: when political insiders lock the channels of democratic change (as Tennessee's rural minority did for 60 years), judicial abstention does not protect democracy---it entrenches minority rule. From this vantage, Justice Frankfurter's appeal to ``an aroused popular conscience'' was circular because malapportionment prevented the majority from translating its conscience into legislative seats (see also \href{https://ir.lawnet.fordham.edu/cgi/viewcontent.cgi?article=2374&context=faculty_scholarship}{\textit{A Politics-Reinforcing Political Question Doctrine}}). Yet even some progressive scholars critique \textit{Baker}'s subsequent trajectory: by channeling voting-structure claims exclusively through individual Equal Protection rights while leaving the Guaranty Clause moribund, \textit{Baker} set the stage for \href{https://supreme.justia.com/cases/federal/us/588/18-422/}{\textit{Rucho v. Common Cause}}, where the Court used \textit{Baker}'s ``judicially manageable standards'' prong to insulate extreme partisan gerrymandering from federal judicial review.

\section{Key Quotations}

\begin{quote}
``Prominent on the surface of any case held to involve a political question is found a textually demonstrable constitutional commitment of the issue to a coordinate political department; or a lack of judicially discoverable and manageable standards for resolving it; or the impossibility of deciding without an initial policy determination of a kind clearly for non judicial discretion; or the impossibility of a court's undertaking independent resolution without expressing lack of the respect due coordinate branches of government; or an unusual need for unquestioning adherence to a political decision already made; or the potentiality of embarrassment from multifarious pronouncements by various departments on one question.''
\end{quote}
\noindent --- Justice Brennan, Opinion of the Court, \href{https://supreme.justia.com/cases/federal/us/369/186/}{369 U.S. at 217}.

\begin{quote}
``The doctrine of which we treat is one of `political questions,' not one of `political cases.' The courts cannot reject as `no law suit' a bona fide controversy as to whether some action denominated `political' exceeds constitutional authority.''
\end{quote}
\noindent --- Justice Brennan, Opinion of the Court, \href{https://supreme.justia.com/cases/federal/us/369/186/}{369 U.S. at 217}.

\begin{quote}
``Have the appellants alleged such a personal stake in the outcome of the controversy as to assure that concrete adverseness which sharpens the presentation of issues upon which the court so largely depends for illumination of difficult constitutional questions? This is the gist of the question of standing.''
\end{quote}
\noindent --- Justice Brennan, Opinion of the Court, \href{https://supreme.justia.com/cases/federal/us/369/186/}{369 U.S. at 204}.

\begin{quote}
``[W]hen the entire table is examined --- comparing the voting strength of counties of like population as well as contrasting that of the smaller with the larger counties --- it leaves but one conclusion, namely that Tennessee's apportionment is a crazy quilt without rational basis.''
\end{quote}
\noindent --- Justice Clark, Concurring Opinion, \href{https://supreme.justia.com/cases/federal/us/369/186/}{369 U.S. at 254}.

\begin{quote}
``The Court's authority --- possessed of neither the purse nor the sword --- ultimately rests on sustained public confidence in its moral sanction. Such feeling must be nourished by the Court's complete detachment, in fact and in appearance, from political entanglements and by abstention from injecting itself into the clash of political forces in political settlements.''
\end{quote}
\noindent --- Justice Frankfurter, Dissenting Opinion, \href{https://supreme.justia.com/cases/federal/us/369/186/}{369 U.S. at 267}.


\end{document}
